\uuid{d63x}
\titre{Équation droite/plan}
\auteur{bodin ; hanani ; mzari}
\organisation{exo7}
\licence{CC BY-NC-SA 4.0}
\datecreate{2020-01-20}
\niveau{L1}
\module{Géométrie}
\chapitre{Géométrie du plan}
\sousChapitre{Distance, norme, produit scalaire}
\difficulte{3}
\theme{équation droite, plan}
\exo7qcmid{lille-1.geometrie-plane.025}

\contenu{
\begin{question}
\qclassification{140.01,140.02}
\qtags{motcle=équation droite/plan}
Dans le plan rapporté à un repère orthonormé, on considère le point $\displaystyle A(1,2)$ et on note $D$ une droite passant par $A$ et qui est à distance $1$ de l'origine. Une équation cartésienne de $D$ est
\begin{answers}
    \good{$D\; :\; x=1$}
    \bad{$D\; :\; x+2y=0$}
    \good{$D\; :\; 3x-4y+5=0$}
    \bad{$D\; :\; y=2x$}
\end{answers}
\begin{explanations}
Une équation cartésienne d'une droite $D$ passant par $A$ est de la forme $a(x-1)+b(y-2)=0$. Mais,
$$1=\mbox{d}(O,D)=\frac{|a+2b|}{\sqrt{a^2+b^2}}\Leftrightarrow b=0\mbox{ ou }b=-\frac{4}{3}a.$$
Ceci détermine toutes les droites passant par $A$ et qui sont à distance $1$ de l'origine.
\end{explanations}
\end{question}
}
